\documentclass[12pt]{article}
\usepackage[a4paper, margin=2cm]{geometry}
\usepackage{amsmath, amssymb, amsthm}
\usepackage{mathtools}
\usepackage{bm}

\theoremstyle{plain}
\newtheorem{theorem}{Theorem}
\newtheorem{lemma}[theorem]{Lemma}

\theoremstyle{definition}
\newtheorem{definition}[theorem]{Definition}

\title{Mathematics Typesetting Test}
\author{TeX Live API}
\date{\today}

\begin{document}
\maketitle

\section{Calculus}

The Fundamental Theorem of Calculus:

\begin{theorem}[Fundamental Theorem of Calculus]
  If $f$ is continuous on $[a,b]$ and $F' = f$, then
  \begin{equation}
    \int_a^b f(x)\,dx = F(b) - F(a). \label{eq:ftc}
  \end{equation}
\end{theorem}

\section{Maxwell's Equations}

\begin{align}
  \nabla \cdot \mathbf{E}  &= \frac{\rho}{\varepsilon_0}           \label{eq:gauss_e} \\
  \nabla \cdot \mathbf{B}  &= 0                                     \label{eq:gauss_b} \\
  \nabla \times \mathbf{E} &= -\frac{\partial \mathbf{B}}{\partial t} \label{eq:faraday} \\
  \nabla \times \mathbf{B} &= \mu_0 \mathbf{J} + \mu_0 \varepsilon_0 \frac{\partial \mathbf{E}}{\partial t} \label{eq:ampere}
\end{align}

\section{Linear Algebra}

A matrix and its determinant:
\begin{equation}
  A = \begin{pmatrix} a_{11} & a_{12} & a_{13} \\ a_{21} & a_{22} & a_{23} \\ a_{31} & a_{32} & a_{33} \end{pmatrix}, \qquad
  \det(A) = \sum_{\sigma \in S_3} \operatorname{sgn}(\sigma) \prod_{i=1}^{3} a_{i,\sigma(i)}
\end{equation}

\section{Piecewise Functions}

\begin{equation}
  f(x) = \begin{cases}
    x^2        & \text{if } x \geq 0 \\
    -x         & \text{if } x < 0 \\
    \text{undefined} & \text{otherwise}
  \end{cases}
\end{equation}

\section{Series and Limits}

Euler's identity and the Basel problem:
\begin{gather}
  e^{i\pi} + 1 = 0 \label{eq:euler} \\[6pt]
  \sum_{n=1}^{\infty} \frac{1}{n^2} = \frac{\pi^2}{6} \label{eq:basel}
\end{gather}

\section{Quantum Mechanics}

The time-independent Schr\"odinger equation:
\begin{equation}
  \hat{H}\Psi = E\Psi, \qquad \hat{H} = -\frac{\hbar^2}{2m}\nabla^2 + V(\mathbf{r})
\end{equation}

\begin{definition}[Hilbert Space]
  A complete inner product space over $\mathbb{C}$ equipped with the norm
  $\|\psi\| = \sqrt{\langle \psi | \psi \rangle}$.
\end{definition}

\begin{proof}
  The completeness follows from the Riesz--Fischer theorem. \qedhere
\end{proof}

\end{document}
